\documentclass[11pt]{article}
\usepackage[T1]{fontenc}
\usepackage{lmodern,amsmath,amssymb,microtype}
\usepackage[margin=1.02in]{geometry}
\usepackage[colorlinks=true,urlcolor=blue,citecolor=blue]{hyperref}
\setlength{\emergencystretch}{2em}
\hypersetup{pdfauthor={Wentao Li},pdftitle={Malicious-verifier statistical zero knowledge for proofs of proximity}}
\title{Malicious-verifier statistical zero knowledge for proofs of proximity}
\author{Wentao Li\\[2pt]\small Prepared with assistance from OpenAI Codex}
\date{November 2026}
\begin{document}
\maketitle
\begin{abstract}
Berman, Rothblum and Vaikuntanathan asked whether the honest-verifier to malicious-verifier transformation for statistical zero knowledge extends to proofs of proximity. We obtain the transformation in their polylogarithmic regime. The verifier and simulator can evaluate the two entropy samplers of the original reduction on chosen seeds using only queries to the large input. This allows the classical public-coin Entropy Difference proof and the strong black-box simulation theorem to be applied with the size of the short oracle circuits as the complexity parameter. The result includes the two graph promises singled out in the original question.
\end{abstract}

\section{Statement}
In a proof of proximity the verifier has query access to an object $f$ of encoded length $N$. Completeness concerns members of the property $\Pi$; soundness concerns objects $\varepsilon$-far from $\Pi$. The prover can read $f$ and is computationally unbounded. The verifier and simulator know the public size parameters, proximity parameter and security parameter $k$. Use the statistical definitions of [1], including auxiliary input for a dishonest verifier. Write $b=1+\lceil\log_2(N+1)\rceil+\lceil1/\varepsilon\rceil$. As in [1], the sizes and addresses in the query representation have efficient binary encodings.

\textbf{Theorem.} Suppose $\Pi$ has an honest-verifier statistical proof of proximity with verifier and honest simulator cost $\operatorname{poly}(b,k)$ in the model of [1, Corollary 5.4]. Then $\Pi$ has a statistical proof of proximity against arbitrary cheating verifiers. Its honest verification time, communication and number of input queries are $\operatorname{poly}(b,k)$. There is a strict black-box simulator of cost $\operatorname{poly}(b,k)(1+t^*)$ up to polynomial simulation overhead for the underlying machine model, where $t^*$ bounds the cheating verifier's time. Completeness error, soundness error and simulation deviation are negligible in $k$. The converse inclusion follows by taking the honest verifier as the verifier strategy.

The quantitative conclusion allows polynomial losses in the original resources. The number of rounds can grow, and polynomial overhead can destroy sublinearity for a verifier initially taking $N^\delta$ time. Polylogarithmic dependence on $N$, at fixed security and proximity, is the regime of the statistical graph problems in [1].

\section{Two ingredients and their interfaces}
\paragraph{Reduction with oracle gates.}
Use [1, full version, Lemma A.1 and Appendix B.1]. On public parameters it constructs two short circuits $C_0,C_1$ with gates querying the input. For uniform seeds let $X=C_0^f(U)$, $Y=C_1^f(U)$. With base-two Shannon entropy,
\[
 f\in\Pi\implies H(X)-H(Y)\ge1,\qquad
 f\text{ is }\varepsilon\text{-far from }\Pi\implies H(Y)-H(X)\ge1 .
\]
The lemma covers expected honest simulation under its stated running-time hypothesis $t'_S(\varepsilon,N)\operatorname{poly}(k)$ with $t'_S$ nondecreasing in $N$. A polynomial bound in $b,k$ admits such a majorant. Its construction time and output size are polynomial in $t_S(\varepsilon,N,\operatorname{poly}(t'_S(\varepsilon,N)))$, hence polynomial in $b$. In particular both circuit descriptions, their seed and output lengths, and the cost of evaluating them on specified seeds are polynomial in $b$. The short descriptions contain oracle gates, never a copy of the table of $f$. The cost of writing an address and reading an answer is included under the query representation of [1]. Duplicate each circuit on two independent seeds if desired: entropy then doubles, so the gap strictly exceeds one in the indicated direction.

\paragraph{Public-coin Entropy Difference with chosen-seed evaluation.}
Goldreich and Vadhan [2, Theorem 1.7, Sections 3--4] give a public-coin honest-verifier proof for the preceding gap promise, with exponentially small simulation deviation and errors. Their verifier and honest simulator work with oracle-aided samplers as well. Here are the input operations that justify this claim.

Their flattening step replaces a distribution by a polynomial tensor power. On an explicit tuple of seeds, this is a list of evaluations of the base sampler. In the main protocol the verifier checks a preimage by evaluating the appropriate tensor sampler on the supplied seed and checks a public hash equation. In the pushing game of Section 4.1 it checks $D(r)=x$ and $h(r,x')=0$. Sample generation and sample testing repeat that game. The honest simulator picks random seeds, evaluates the associated samplers, and draws public hashes subject to specified linear constraints. In particular, when its testing simulator needs coins for a sample, the main simulator generated the sample together with those coins. Any enumerations of seed preimages required of the honest prover are permitted. The verifier and simulator incur only polynomially many evaluations on known seeds and polynomial work on the short strings and hash descriptions. Merely drawing a fresh independent sample on demand would be an insufficient implementation of the verifier's preimage check. Here a specified seed is evaluable.

More formally, fix $f$ and run this algorithm with evaluation subroutines for $C_0^f,C_1^f$. The distribution of every message and the final decision coincide with those of the algorithm on the same two deterministic seed maps. The analysis of [2, Secs.~3--4] uses their input and output lengths, entropy, and hash distributions. Its security parameter $s$ can be chosen independently of any Boolean circuit that would hardwire the full table of $f$. Choose $s$ at least the sum of $k$, the short description lengths and a sufficiently large absolute constant. Equivalently, use a padded public instance of that length with oracle-gate descriptions. This gives a public-coin protocol with polynomial cost in $b+k$ and error exponentially small in $s$, uniformly over the oracle. This access argument is also consistent with the relativization of honest-verifier versus general SZK recorded in [4, Sec.~2]. Relativization alone would leave the cost with respect to $N$ to be checked.

\section{Compiling and simulating}
Apply [3, Theorem 3] to this public-coin protocol, with the public padded size as length parameter. Its preliminary amplification and public-coin random-selection steps are uniform in the verifier and honest simulator. In particular, at each compiled round the parties jointly generate the challenge that the original verifier would have sent. The compiled verification predicate then runs the original public-coin verifier. The strong simulator uses the original honest simulation and black-box calls to the cheating strategy to simulate challenge selection. Hashing and selection in the compiler use short public strings. Accesses to $f$ in the honest verification and its compiled simulator can thus be implemented by the chosen-seed evaluations above and, for simulation, by running the cheating strategy with its original oracle access. Neither step inspects arbitrary entries in an expanded circuit description. One may fix $f$ throughout the security proof. All the estimates in [2,3] are uniform in the seed maps of the given sizes.

By [3, Theorem 3], the strong simulation deviation is at most a polynomial in the padded size times the honest simulation error, plus an exponentially vanishing term. Completeness error is exponentially small. Choose parameters to make soundness at most $1/3$ (Theorem 3 permits inverse-polynomial soundness). The number of ordinary steps, oracle evaluations and black-box activations per execution is polynomial in $b+k$. The simulator is strict polynomial time in the black-box cost convention of [3, Sec.~1.1]. For a cheating strategy with arbitrary auxiliary input $z$, fix its random tape across each rewind and answer its own oracle queries by the same $f$. A black-box activation costs at most $O(1+t^*)$ up to ordinary emulation overhead; coins on the portion of tape accessed can be generated lazily with consistent replay. The simulator need only present to the strategy the same access to $z$ it has in the real experiment. Its runtime is independent of the total length of unread auxiliary data. The strong guarantee in [3] is uniform in the cheating strategy, including its auxiliary state.

Finally perform $k+1$ independent sequential executions, accepting if every execution accepts. On a far input, conditional on a prior transcript, a prover's continuation is still a strategy in a fresh execution, so soundness is at most $(1/3)^{k+1}$. On a member, completeness error is at most the sum of the single-execution errors. For simulation, proceed sequentially while retaining the cheating verifier's state and its consistent random tape. At each step regard its saved state as auxiliary input for the continuation. The universal auxiliary-input guarantee and a hybrid over executions bound the total deviation by $(k+1)$ times the single-execution deviation. Equivalently one may invoke sequential composition as in [3, Sec.~1.3]. Protocol time and the cost of storing or replaying the interaction stay polynomial in $b+k$ times the cost of running the verifier. Since $s\ge k$, the error estimates are negligible in $k$. This proves the theorem.

\section{Application and provenance}
Apply the theorem to the bounded-degree promise expansion and promise bipartiteness protocols of [1, ITCS Theorem 3; full version, Secs.~4.2--4.3]. These now protect against arbitrary verifier strategies with polynomial cost in $\log N,k,1/\varepsilon$ under the original graph hypotheses. For example the bipartiteness soundness promise includes rapid mixing on the far side. Removing that hypothesis is a different question. These are interactive, information-theoretic proofs in the input-query model.

[1] I. Berman, R. D. Rothblum and V. Vaikuntanathan, \href{https://eprint.iacr.org/2017/114}{Zero-Knowledge Proofs of Proximity}, 2017 preprint, footnote 7, Lemma A.1 and Corollary 5.4. \href{https://doi.org/10.4230/LIPIcs.ITCS.2018.19}{ITCS 2018}, footnote 10.\par
[2] O. Goldreich and S. Vadhan, \href{https://doi.org/10.1109/CCC.1999.766262}{Comparing Entropies in Statistical Zero Knowledge with Applications to the Structure of SZK}, CCC 1999; \href{https://www.wisdom.weizmann.ac.il/~oded/PSX/gv2.pdf}{full paper}, Theorem 1.7, Corollary 1.9, Sections 3--4.\par
[3] O. Goldreich, A. Sahai and S. Vadhan, \href{https://people.seas.harvard.edu/~salil/research/dishonest-abs.html}{Honest-Verifier Statistical Zero-Knowledge Equals General Statistical Zero-Knowledge}, STOC 1998, Theorem 3 and Section 1.3.\par
[4] A. Bouland, L. Chen, D. Holden, J. Thaler and P. N. Vasudevan, \href{https://arxiv.org/abs/1609.02888}{On the Power of Statistical Zero Knowledge}, Sec.~2. They explicitly note that the classical equivalence holds with an oracle. The contribution of this note is the composition under the proximity resource bounds. The search located no prior explicit resolution of the proximity question. Independent priority and proof review are pending.
\end{document}
